% Part:sets-functions-relations
% Chapter: sets
% Section: reduction

\documentclass[../../../include/open-logic-section]{subfiles}

\begin{document}

\olfileid[pt]{sfr}{siz}{red}

\olsection{Redução}

\begin{editorial}
  Esta seção prova a não enumerabilidade por redução, correspondendo aos
  resultados na \olref[nen]{sec}. Uma versão alternativa, ligeiramente mais condensada,
  correspondente aos resultados na \olref[nen-alt]{sec}, é fornecida na
  \olref[red-alt]{sec}.
\end{editorial}

Mostramos $\Pow{\PosInt}$ como !!{nonenumerable} por um argumento de diagonalização.
Já tínhamos uma prova de que $\Bin^\omega$, o conjunto de todas as
sequências infinitas de $0$s e $1$s, é !!{nonenumerable}. Eis
outra maneira de provar que $\Pow{\PosInt}$ é !!{nonenumerable}:
mostre que \emph{se $\Pow{\PosInt}$ é !!{enumerable}, então $\Bin^\omega$
  também é !!{enumerable}}. Como sabemos que $\Bin^\omega$ não é
!!{enumerable}, $\Pow{\PosInt}$ também não pode ser. Isso é chamado de
\emph{reduzir} um problema a outro---neste caso, reduzimos o
problema de enumerar $\Bin^\omega$ ao problema de enumerar
$\Pow{\PosInt}$. Uma solução para o último---uma enumeração de
$\Pow{\PosInt}$---daria uma solução para o primeiro---uma enumeração
de $\Bin^\omega$.

Como reduzimos o problema de enumerar um conjunto~$B$ ao de
enumerar um conjunto~$A$? Fornecemos uma maneira de transformar uma enumeração
de~$A$ em uma enumeração de~$B$. A maneira mais fácil de fazer isso é
definir uma função !!{surjective} $f\colon A \to B$. Se $x_1$, $x_2$,
\dots{} enumera~$A$, então $f(x_1)$, $f(x_2)$, \dots{} enumeraria~$B$. Em nosso caso, procuramos uma função sobrejetora
$f\colon \Pow{\PosInt} \to \Bin^\omega$.

\begin{prob}
Mostre que se há uma função !!{injective} $g\colon B \to A$, e
$B$~é !!{nonenumerable}, então $A$ também é. Faça isso mostrando como você
pode usar~$g$ para transformar uma enumeração de~$A$ em uma de~$B$.
\end{prob}

\begin{proof}[Prova do {\olref[nen]{thm:nonenum-pownat}} por redução]
Suponha que $\Pow{\PosInt}$ fosse !!{enumerable}, e portanto que
há uma enumeração dele, $Z_{1}$, $Z_{2}$, $Z_{3}$, \dots

Defina a função $f \colon \Pow{\PosInt} \to \Bin^\omega$ fazendo
$f(Z)$ ser a sequência $s_{k}$ tal que $s_{k}(n) = 1$ sse $n \in Z$,
e $s_k(n) = 0$ caso contrário. Isso define claramente uma função, já que
sempre que $Z \subseteq \PosInt$, qualquer $n \in \PosInt$ ou é
!!a{element} de $Z$ ou não é. Por exemplo, o conjunto $2\PosInt = \{2,
4, 6, \dots\}$ dos pares positivos é mapeado para a sequência
$010101\dots$, o conjunto vazio é mapeado para $0000\dots$ e o conjunto
$\PosInt$ para $1111\dots$.

Também é !!{surjective}: toda sequência de $0$s e $1$s corresponde
a algum conjunto de inteiros positivos, nomeadamente aquele que tem como
membros os inteiros correspondentes às posições onde a sequência
tem~$1$s. Mais precisamente, suponha $s \in \Bin^\omega$. Defina $Z
\subseteq \PosInt$ por:
\[
Z = \Setabs{n \in \PosInt}{s(n) = 1}
\]
Então $f(Z) = s$, como pode ser verificado consultando a definição
de~$f$.

Agora considere a lista
\[
f(Z_1), f(Z_2), f(Z_3), \dots
\]
Como $f$ é !!{surjective}, todo membro de $\Bin^\omega$ deve
aparecer como valor de~$f$ para algum argumento, e assim deve aparecer na
lista. Essa lista portanto deve enumerar todo~$\Bin^\omega$.

Assim, se $\Pow{\PosInt}$ fosse !!{enumerable}, $\Bin^\omega$ seria
!!{enumerable}. Mas $\Bin^\omega$ é !!{nonenumerable}
(\olref[nen]{thm:nonenum-bin-omega}). Logo $\Pow{\PosInt}$ é
!!{nonenumerable}.
\end{proof}

\begin{explain}
É fácil confundir a direção em que a redução vai.
Por exemplo, uma função !!{surjective} $g \colon \Bin^\omega \to B$
\emph{não} estabelece que $B$ é !!{nonenumerable}. (Considere $g
\colon \Bin^\omega \to \Bin$ definida por $g(s) = s(1)$, a função
que mapeia uma sequência de $0$s e $1$s para seu primeiro !!{element}. Ela
é !!{surjective}, porque algumas sequências começam com $0$ e outras começam
com $1$. Mas $\Bin$ é finito.) Note também que a função~$f$ deve
ser !!{surjective}, caso contrário o argumento não se sustenta:
$f(x_1)$, $f(x_2)$, \dots{} não estariam garantidos de incluir
todos os !!{element}s de~$B$. Por exemplo,
\[
h(n) = \underbrace{000\dots0}_{\text{$n$ $0$s}}
\]
define uma função $h\colon \PosInt \to
\Bin^\omega$, mas $\PosInt$ é !!{enumerable}.
\end{explain}

\begin{prob}
Mostre que o conjunto de todos os \emph{conjuntos de} pares de inteiros positivos é
!!{nonenumerable} por um argumento de redução.
\end{prob}

\begin{prob}\label{sfr:siz:red:prob:nat-nat}
  Mostre que o conjunto~$X$ de todas as funções $f\colon \Nat \to \Nat$ é
  !!{nonenumerable} por um argumento de redução (Dica: dê uma função sobrejetora
  de $X$ para~$\Bin^\omega$.)
\end{prob}

\begin{prob}
Mostre que $\Nat^\omega$, o conjunto de sequências infinitas de
números naturais, é !!{nonenumerable} por um argumento de redução.
\end{prob}

\begin{prob}
Seja $P$ o conjunto de funções do conjunto dos inteiros positivos
para o conjunto $\{0\}$, e seja $Q$ o conjunto de funções \emph{parciais}
do conjunto dos inteiros positivos para o conjunto $\{0\}$. Mostre
que $P$~é !!{enumerable} e $Q$~não é. (Dica: reduza o problema
de enumerar $\Bin^\omega$ a enumerar~$Q$).
\end{prob}

\begin{prob}
Seja $S$ o conjunto de todas as funções !!{surjective}s do conjunto dos
inteiros positivos para o conjunto \{0,1\}, isto é, $S$ consiste de todas as
!!{surjective}s~$f\colon \PosInt \to \Bin$. Mostre que $S$ é
!!{nonenumerable}.
\end{prob}

\begin{prob}
Mostre que o conjunto~$\Real$ de todos os números reais é !!{nonenumerable}.
\end{prob}

\end{document}
